\documentclass[../main.tex]{subfiles}

\begin{document}

    We were given the equation $f(x) = (x + 5)^{1/3} - x + 3 = 0$ in the interval $x \in [3,7]$ and must use
    the fixed-point method to find a root, $x_{*}$, for any $g(x)$ that converges and verify
    the root has been found within a numerical error - plotting our results.

    Our $g(x)$ functions (iteration functions) are defined as:
    \begin{itemize}
        \item $g_{1}(x) = (x + 5)^{1/3} + 3$,
        \item $g_{2}(x) = (x - 3)^{3} - 5$.
    \end{itemize}
    
    \section{Convergence of iteration functions}\label{sec:fixed_point_convergence}
    For some $g_{n}(x)$ to converge, the following conditions must be met:
    \begin{itemize}
        \item $f(x)$ must be continuous on the closed interval $[3,7]$,
        \item $f(x)$ must be differentiable on the open interval $(3,7)$,
        \item there exists $\xi \in (3,7)$ that brackets the root,
        \item and some $|g_{n}'(\xi)| < 1$.
    \end{itemize}

    \subsection*{Is $f(x)$ continuous on $[3,7]$?}\label{subsec:fixed-point-continuous}
    \underline{Proof:} $(x + 5)$ and $(-x + 3)$ are polynomials, so they are continuous everywhere.
    The function $a(x) = x^{1/3}$ is continuous $\forall \ x \in \mathbb{R}$,
    hence $a(x + 5)$ is continuous $\forall \ x \in \mathbb{R}$.
    So, $a(x + 5) + (-x + 3)$ is continuous $\forall \ x \in \mathbb{R}$.
    Therefore, $f(x)$ is continuous on the closed interval $[3,7]$.

    \subsection*{Is $f(x)$ differentiable on $(3,7)$?}\label{subsec:fixed-point-differentiable}
    \underline{Proof:} $(x + 5)$ and $(-x + 3)$ are polynomials, so they are differentiable everywhere.
    The function $a(x) = x^{1/3}$ is differentiable $\forall \ x \in \mathbb{R}$ with $x \neq 0$.
    $x \in (3,7)$, so $x + 5 \in (8,12)$ and therefore $x + 5 \neq 0$.
    Hence, $a(x + 5)$ is differentiable $\forall \ x \in \mathbb{R}$.
    So, $a(x + 5) + (-x + 3)$ is differentiable $\forall \ x \in \mathbb{R}$.
    Therefore, $f(x)$ is differentiable on the open interval $(3,7)$.

    \subsection*{Does $(3,7)$ bracket a root?}\label{subsec:fixed-point-brackets-root}
    This is simple enough to prove.
    The signs of the values of $f(x)$ must be opposite for the minimum and maximum of the interval $(3,7)$.
    \begin{itemize}
        \item $f(3) = (3 + 5)^{1/3} - 3 + 3 = 8^{1/3} = +2$
        \item $f(7) = (7 + 5)^{1/3} - 7 + 3 = 12^{1/3} - 4 \approx -1.71$
    \end{itemize}\\
    The two ends of our interval have opposite signs, hence there is a root in $(3,7)$.

    \subsection*{Do the iterable functions converge?}\label{subsec:fixed-point-iterables-convergence}
    For a function $g_{n}(x)$ to converge $|g_{n}'(x)| < 1$, so:
    \begin{itemize}
        \item Does $g_{1}(x) = (x_{n} + 5)^{1/3} + 3$ converge? \\
        \\ \underline{Proof:} \\
        \[g_{1}'(x) = \frac{1}{3}(x_{n} + 5)^{-2/3}\] \\
        $(x + 5)^{-2/3}$ is decreasing for $x > 0$, so its maximum on $[3,7]$ occurs at $x = 3$.
        \[\max_{x \in (3,7)} g_{1}(x) = \frac{1}{3}(3 + 5)^{-2/3} =
        \frac{1}{3} \cdot 8^{-2/3} = \frac{1}{3} \cdot \frac{1}{4} = \frac{1}{12}\] \\
        So, $\forall \ \xi \in (3,7)$ we have that $|g_{1}'(\xi)| < 1$ and $g_{1}(x)$ converges.

        \item Does $g_{2}(x) = (x_{n} - 3)^{3} - 5$ converge? \\
        \\ \underline{Proof:} \\
        \[g_{2}'(x) = 3(x_{n} - 3)^{2}\] \\
        $(x - 3)^{2}$ is increasing for $x > 0$, so its maximum on $[3,7]$ occurs at $x = 7$.
        \[\max_{x \in (3,7)} g_{2}(x) = 3(7 - 3)^{2} = 3(4)^{2} = 3 \cdot 16 = 48 \] \\
        So, $\exists \ \xi \in (3,7)$ where $|g_{2}'(x)| \geq 1$, hence $g_{2}(x)$ does not converge.
    \end{itemize}

    \subsection*{Conclusion}\label{subsec:fixed-point-conclusion}
    We have found that a root exists in $[3,7]$ and that $g_{1}(x)$ can be used to find it, as it will converge to $x_{*}$.

    \section{Fixed-point iteration}\label{sec:fixed_point_iteration}

    \subsection*{Proving that $g_{2}(x)$ diverges}\label{subsec:fixed-point-prove-g2-diverges}
    First, let's quickly look at $g_{2}(x)$:
    \begin{center*}
        \begin{align*}
            & x_{n+1} = g_{2}(x_{n}) = (x_{n} - 3)^{3} - 5 && \\
            & x_{0} = 6 && \\
            & x_{1} = g_{2}(6) = (6 - 3)^{3} - 5 = 3^{3} -5 = 22 && \\
            & x_{2} = 6854 && \\
            & x_{3} \approx \num{3.2156e11} && \\
            & \vdots &&
        \end{align*}
    \end{center*}\\
    We can see that, as $n$ increases, $x_{n}$ exponentially increases
    showing numerically that $g_{2}(x)$ does not converge to $x_{*}$.

    \subsection*{Finding the root with $g_{1}(x)$}\label{subsec:fixed-point-finding-root}
    Now, let's look at $g_{1}(x)$:
    \begin{center*}
        \begin{align*}
            & x_{n+1} = g_{1}(x_{n}) = (x_{n} + 5)^{1/3} + 3 && \\
            & x_{0} = 6 && \\
            & x_{1} = g_{1}(6) = (6 + 5)^{1/3} + 3 = 11^{1/3} + 3 \approx 5.2240 && \\
            & x_{2} \approx 5.1704 && \\
            & x_{3} \approx 5.1666 && \\
            & \vdots && \\
            & x_{10} \approx 5.1663 &&
        \end{align*}
    \end{center*}\\
    We can see that as $n$ increases, $x \to x_{*}$.

    \subsection*{Verifying the root}\label{subsec:fixed-point-verify-root}
    Before we get too excited, let's check that what we have found is our root:
    \begin{center*}
        \begin{align*}
            & x_{*} = x_{10} = 5.1663 && \\
            & f(x_{*}) = (x_{*} + 5)^{1/3} - x + 3 \approx \num{-2.46e-12} \approx 0 && \\
            & |f(x_{10}) - f(x_{9})| = \num{-3.4637e-11} - \num{-4.8765e-10} \approx \num{4.5301e-10} \approx 0 &&
        \end{align*}
    \end{center*}

    \subsection*{Code used to find the root}\label{subsec:fixed-point-code}
    I ran the following in my terminal (usage covered in
    \hyperref[subsubsec:fixed-point-method-appendix]{Fixed-point method} under
    \hyperref[subsec:usage-of-modules]{Usage of modules}):
    \begin{lstlisting}[style=console,label={lst:fixed-point-method}]
C:\User\CaseStudy> python main.py fixed-point ./images
Root: 5.1663
Error: 0.0000
    \end{lstlisting}

    Producing these graphs:

    \begin{figure}[H]
        \centering
        \begin{subfigure}
            \centering
            \includegraphics[width=230px]{fixed-point-all-functions}
            \label{fig:fixed-point-all-functions}
        \end{subfigure} \\
        \begin{subfigure}
            \includegraphics[width=230px]{fixed-point-g1}
            \label{fig:fixed-point-g1-iterations}
        \end{subfigure}
        \begin{subfigure}
            \includegraphics[width=230px]{fixed-point-g2}
            \label{fig:fixed-point-g2-iterations}
        \end{subfigure}
        \textbf{Figure:} $f(x)$, $g_{1}(x)$, $g_{2}(x)$, $y = x$, and $g_{n}(x)$ iterations plotted.
        \label{fig:fixed-point-plots}
    \end{figure}

\end{document}
